 % 本文件由 scripts/assemble.py 自动装配，勿手改（改 parts/ 后重跑）
% 第1章 引言

% 第1章 INTRODUCTION

\chapter{引言}

\section{多者异也}

\begin{keypoints}
\item 多体系统的集体激发可以看作粒子。然而，这些粒子的性质可能与构成多体系统的粒子的性质非常不同。
\item 猜想胜于推导。
\item 经典计算的极限。
\item 我们的真空不过是一种特殊的材料。
\end{keypoints}

\emph{量变可以导致质变。}这一哲理在含有大量粒子（或大量自由度）的系统（如固体和液体）中一遍又一遍地得到演示。支配少数粒子系统的物理原理，可以与支配多体系统集体运动的物理原理非常不同。新的物理概念（例如费米子和规范玻色子的概念）以及新的物理定律和原理（例如电磁学定律），可以从大量粒子的关联中产生（见第 10 章）。

凝聚态物理是物理学的一个分支，研究处于“凝聚”（即固体或液体）状态的多粒子系统。当前凝聚态理论的出发点是支配许多粒子（例如电子和原子核）运动的薛定谔方程。薛定谔方程在数学上是完备的。原则上，我们可以通过求解相应的薛定谔方程，得到任何一个多体系统的全部性质。

然而在实践中，所需的计算能力大得惊人。20 世纪 80 年代，一台具有 32 Mbyte RAM 的工作站可以求解十一个相互作用电子的系统。二十年之后，计算能力提高了 100 倍，这不过让我们多求解了两个电子。求解一个典型的、含 $10^{23}$ 个相互作用电子的系统所需的计算能力，超出了人类大脑的想象。即使用我们宇宙中的全部原子造出一台经典计算机，也不足以处理这个问题\footnote{它甚至连存储该系统单个态矢量的内存都不够。}。这样一台不可能的计算机也只能求解大约 100 个粒子的薛定谔方程\footnote{这引出一个非常有意思的问题——自然界是如何做计算的？自然界是如何算出 $10^{23}$ 个粒子一秒之后的状态的？看来我们所用的数学太低效了。自然界并不这样计算。}。我们看到，一个一般的相互作用多体系统是一个极端复杂的系统。实际上，不可能从薛定谔方程推导出它的全部确切性质。所以，即使薛定谔方程是凝聚态系统的正确理论，它也未必总是有助于获得相互作用多体系统的物理性质。

即便我们真的得到了一个一般相互作用多体系统的精确解，结果也往往复杂到几乎不可能被完全透彻地理解。为了体会这一结果的复杂程度，让我们考虑一个仅含 200 个电子的小型相互作用系统。系统的能量本征值分布在约 $200\,\mathrm{eV}$ 的范围内。该系统至少有 $2^{200}$ 个能级，能级间隔约为 $200\,\mathrm{eV}/2^{200}=10^{-60}\,\mathrm{eV}$。假如我们花费相当于宇宙年龄的时间来测量能量，那么由于能量–时间不确定性关系，我们也只能达到 $10^{-33}\,\mathrm{eV}$ 量级的能量分辨率。可见，相互作用多体系统的精确结果可以复杂到无法在实验上对其正确性作完全透彻的检验\footnote{由于我们无法完全透彻地检验由薛定谔方程得到的结果，“薛定谔方程决定多体系统的所有性质”这一信念只不过是一种\emph{信仰}。}。要真正理解一个系统，我们需要理解系统各种不同现象之间的联系与关系，而薛定谔方程往往并不能直接提供这种理解。

由于我们通常无法直接用薛定谔方程来理解相互作用系统，当我们面对一个多体系统时，就必须从头开始。我们必须把多体系统当作一个黑箱来处理，就像我们对待神秘而未知的宇宙那样。我们必须去\emph{猜想}一个直接联系各种实验观测的低能有效理论，而不是从薛定谔方程把它推导出来。我们不能想当然地认为，描述低能激发的理论与描述底层电子和原子核的理论有任何相似之处。

这种思路与高能物理的思路非常相似。确实，强关联多体系统的研究与高能物理的研究有着根深蒂固的相似之处：在两种情形下，人们都试图寻找能把一个观测到的实验事实与另一个联系起来的理论。（其实，把一个观测到的实验事实与另一个联系起来，几乎就是物理理论的定义。）一个主要差别是：在高能物理中我们只有一种“材料”（我们的真空）可供研究，而在凝聚态物理中有许多不同的材料，其中可能包含我们的真空中没有的新现象（例如分数统计、非阿贝尔统计，以及具有各种各样规范群的规范理论）。

\section{“基本”粒子与物理定律是演生现象}

\begin{keypoints}
\item 演生——多体系统的第一原理。
\item “基本”粒子的起源。
\item 物理定律之“美”的起源。（为什么自然界表现得合情合理。）
\end{keypoints}

历史上，在探索理解自然的道路上，我们一直被一个基本的（而且不正确的）假设所误导，即真空是空的。我们（不正确地）假设，放在真空中的物质总可以被分割成更小的部分。我们不断地把物质分割成越来越小的部分，试图发现最小的“基本”粒子——我们宇宙的基本砌块。我们一直相信，支配“基本”粒子的物理定律必定是简单的；自然界中丰富多彩的现象都来自这些简单的物理定律。

然而，多体系统呈现出一张非常不同的图景。在高能量（或高温度）和短距离下，多体系统的性质由形成系统的原子/分子之间的相互作用所控制，这种相互作用可以非常复杂而且特殊。当我们降低温度时，视原子之间相互作用形式的不同，会形成晶体结构或超流态。在晶体或超流体中，唯一的低能激发是原子的集体运动，这些激发就是声波。在量子理论中，所有的波都对应于粒子，与声波对应的粒子叫做声子（phonon）\footnote{晶体有三种声子，而超流体只有一种声子。}。因此，在低温下，一个由新型粒子——声子——所支配的新“世界”演生出来了。声子的世界是一个简单而“美丽”的世界，与原来的原子/分子系统大不相同。

让我们解释一下“声子的世界是简单而美丽的”是什么意思。为简单起见，我们将集中考虑超流体。虽然气体中原子之间的相互作用可以复杂而特殊，但低能下演生声子的性质却是简单而普适的。例如，不论原子之间相互作用的形式如何，所有声子都具有与能量无关的速度；尽管原子之间有强相互作用，声子彼此穿过时却几乎不发生相互作用。除声子之外，超流体还有另一种叫做旋子（roton）的激发。旋子之间可以通过交换声子发生相互作用，这导致一种力正比于 $1/r^4$ 的偶极相互作用。我们看到，不仅声子是演生的，就连支配声子和旋子低能世界的物理定律也是演生的。演生的物理定律（例如偶极相互作用定律和无相互作用声子定律）是简单而美丽的。

我认为 $1/r^4$ 偶极相互作用定律是美丽的，因为它不是 $1/r^3$，不是 $1/r^{4.13}$，也不是亿万种其他可能选择之一。它恰恰就是 $1/r^4$，因此，理解它为什么必须是 $1/r^4$ 就十分令人着迷。类似地，$1/r^2$ 库仑定律同样美丽而令人着迷。我们将在本书的前半部分解释超流体中偶极相互作用定律的演生，在本书的后半部分解释库仑定律的演生。

如果我们的宇宙本身是一个超流体，而形成超流体的粒子尚待发现，那么我们将只知道低能声子。我们很可能禁不住要把声子当作基本粒子，把旋子之间的 $1/r^4$ 偶极相互作用当作自然界的基本定律。很难想象，这些声子以及 $1/r^4$ 偶极相互作用定律，竟然来自受一套完全不同的定律支配的粒子。

我们看到，在多体系统中，支配演生低能集体激发的定律是简单的，而且这些集体激发的行为像粒子。如果我们要在多体系统与我们的真空之间建立联系，那么就应当把多体系统中的低能集体激发与真空中的“基本”粒子（例如光子和电子）联系起来。但是，在多体系统中，集体激发并不是基本的。当我们在短长度尺度上审视它们时，看到的是一个复杂的、非普适的原子/分子系统。因此，在多体系统中，我们在低能下拥有集体激发（也称为准粒子），而这些集体激发往往并不成为高能量、短距离下模型的砌块。原子尺度上的理论通常是复杂、特殊而且不合情理的。支配集体激发的物理定律，其简单性和美丽并非来自原子/分子模型的简单性，而是来自这样一个事实：这些定律必须允许集体激发在低能下存活。集体激发之间的一般的相互作用可能给这些激发一个大的能隙，使它们在低能下无法被观察到。能够允许无能隙（或几乎无能隙）集体激发存在的相互作用（或物理定律）必定非常特殊——而且“美丽”。

如果我们相信真空可以被看作一种特殊的多体材料，那么我们就不得不得出结论：不存在“基本”粒子。真空中所有所谓的“基本”粒子实际上都是低能集体激发，它们未必是基本理论的砌块。高能量\footnote{这里，“高能量”指普朗克量级的能量 $M_P=1.2\times10^{19}$ GeV。}、短距离下的基本理论及其砌块，受另一套物理定律支配。从演生的观点看，那些定律可能是特殊的、非普适的、复杂的。

低能量、长距离下美丽的世界与合情合理的物理定律，是作为一种“自然选择”的结果演生出来的：支配低能激发的物理定律，必须允许这些激发在低能下存在。在某种意义上，这种“自然选择”解释了为什么我们的世界是合情合理的。

既懂凝聚态物理又懂高能物理的人可能会反对上述图景，因为我们的真空似乎与我们熟悉的固体和液体非常不同。例如，我们的真空包含狄拉克费米子（例如电子和夸克）和规范玻色子（例如光），而固体和液体看起来并不包含这类激发。看来光和电子是基本的，不可能是演生的。所以，要把多体系统中的演生图景应用于基本粒子，我们必须回答下面这个问题：规范玻色子和狄拉克费米子能从多体系统中演生出来吗？或者，更有意思的是，规范玻色子和狄拉克费米子能从多玻色子系统中演生出来吗？

这里的根本问题是：费米子和规范玻色子从何而来？光与费米子的起源是什么？光和费米子可能是演生现象吗？我们知道，无质量（即无能隙）的粒子在自然界中非常罕见。如果它们存在，那么它们的存在必定有其原因。但是，无质量光子和近无质量费米子（例如电子）得以存在的背后原因是什么？（电子质量比自然标度——普朗克质量——小 $10^{22}$ 倍，对我们的目的而言可以视为零。）多体系统能为上述问题提供答案吗？

在下面几节中，我们将讨论多体系统的一些基本概念。特别是，我们将讨论导致无能隙激发的概念，以及从局域玻色模型导出演生规范玻色子和费米子的概念。我们将看到，无质量光子和无质量费米子可以是演生现象。

\section{凝聚态物理的基石}

\begin{keypoints}
\item 朗道对称性破缺理论（加上重整化群理论）和朗道费米液体理论构成了传统凝聚态物理的基础。
\end{keypoints}

传统多体理论建立在两块基石之上，即朗道费米液体理论和朗道对称性破缺理论（Landau, 1937; Ginzburg and Landau, 1950）。费米液体理论是围绕某一类特殊基态——用单粒子能级填充所得到的态——展开的微扰论。它描述金属、半导体、磁体、超导体和超流体。朗道对称性破缺理论指出，不同的相之所以不同，是因为它们具有不同的对称性；相变不过是改变对称性的转变。朗道对称性破缺理论描述了几乎所有已知的相，例如固相、铁磁相与反铁磁相、超流相等，以及它们之间所有的相变。

我们先撇开光与费米子的起源，来考虑一个较简单的问题——声子的起源。利用朗道对称性破缺理论，我们可以理解无能隙声子的起源。在朗道对称性破缺理论中，如果系统的基态具有一种叫做连续对称性自发破缺的特殊性质（Nambu, 1960; Goldstone, 1961），则该相可以拥有无能隙激发。固体中存在无能隙声子，是因为固体破缺了连续平移对称；无能隙声子恰好有三种，则是因为固体破缺了 $x$、$y$、$z$ 三个方向的平移对称。因此，我们可以说，无能隙声子的起源是固体中的平移对称性破缺。

十分有意思的是，我们对一种无能隙激发——声子——的理解，植根于我们对物质相的理解。知道了光是一种无质量激发之后，人们或许会猜想：光是否也和声子一样，是来自某个破缺对称的 Nambu–Goldstone 模（Nambu–Goldstone mode）？然而，实验告诉我们，像光这样的规范玻色子，在 $3+1$ 维中确实不同于 Nambu–Goldstone 模。

20 世纪 70 年代末，我们觉得至少在原则上已经理解了关于相与相变的全部物理。在朗道对称性破缺理论中，如果我们从纯玻色模型出发，那么得到无能隙激发的唯一途径是连续对称性的自发破缺，而这将导致无能隙的\emph{标量玻色}激发。看来没有办法从对称性破缺得到无能隙规范玻色子和无能隙费米子。这也许就是人们认为我们的真空（具有无质量规范玻色子和近无能隙费米子）与玻色多体系统（当时认为只包含无能隙标量玻色集体激发，例如声子）非常不同的原因。似乎不存在任何序能给出无质量光和无质量费米子。因此，我们把光和费米子归入与声子不同的范畴：我们把它们看作基本的，并人为地引进到我们的自然理论之中。

然而，如果我们真的相信光和费米子像声子一样是事出有因的，那么这个原因必定是真空中某种保护它们无质量性的序\footnote{这里我们已经假定，光和费米子并不是我们放进空真空里的东西。我们的真空更像一个并不空虚的“海洋”。光和费米子是对应于“水”运动之特定模式的集体激发。}。现在的问题是：什么样的序能产生光和费米子，并保护它们的无质量性？从这个观点看，光和费米子的存在本身就表明，我们对物质状态的理解是不完备的。我们应当深化并拓展对物质状态的理解。应当存在包含新型序的新物质状态。新的序将产生光和费米子，并保护它们的无质量性。

\section{拓扑序与量子序}

\begin{keypoints}
\item 在朗道理论之外还有一个新世界。这个新世界丰富多彩、激动人心。
\end{keypoints}

我们对这种新序的理解始于一个意想不到的地方——分数量子霍尔（fractional quantum Hall，FQH）系统。1982 年发现的 FQH 态（Tsui et al., 1982; Laughlin, 1983）揭开了凝聚态物理的新篇章。FQH 态真正新颖之处在于：我们失去了传统多体理论的两块基石。由于量子霍尔系统中很强的相互作用与关联，朗道费米液体理论不适用于这些系统。更为惊人的是，FQH 系统在零温下包含许多\emph{不同}的相，而这些相却具有\emph{相同}的对称性。因此，这些相无法用对称性来区分，也无法用朗道对称性破缺理论来描述。我们突然发现，传统多体理论中没有任何东西可以用来处理这些新问题。于是，强关联系统领域的理论进展要求引入新的数学技巧和物理概念，以超越费米液体理论和朗道对称性破缺原理。

在强关联系统领域，高能粒子理论的发展与凝聚态理论的发展确实相互滋养。我们已经看到许多场论技巧——例如非线性 $\sigma$ 模型、规范理论、玻色化、流代数等——被引进到强关联系统和无序系统的研究之中。这带来了该领域的迅猛发展，以及超越费米液体理论和朗道对称性破缺理论的新理论。本书试图介绍凝聚态理论中这些新进展的一部分。

新进展之一是量子序/拓扑序的引入。由于 FQH 态无法用朗道对称性破缺理论描述，有人提出 FQH 态包含一种新型的序——拓扑序（topological order；Wen, 1990, 1995）。拓扑序之所以新，是因为它不能用对称性破缺、长程关联或局域序参量来描述；我们用来刻画一个相的所有惯用工具都不适用于拓扑序。尽管如此，拓扑序并不是一个空洞的概念，因为它可以用一整套新工具来刻画，例如简并基态的数目（Haldane and Rezayi, 1985）、准粒子统计（Arovas et al., 1984）以及边缘态（Halperin, 1982; Wen, 1992）。

\begin{figure}[H]
\centering
\includegraphics[width=0.72\textwidth]{fig_1.1.pdf}
\caption*{图 1.1\quad 物质中（以及真空中）不同序的分类。}
\end{figure}

已经证明，拓扑有序态的基态简并对\emph{任何}微扰都是稳健的（Wen and Niu, 1990）。因此，基态简并是一种可以用来刻画相的普适性质。拓扑简并基态的存在证明了拓扑序的存在。拓扑简并还可以用来执行容错量子计算（Kitaev, 2003）。

拓扑序的概念最近被推广为量子序（quantum order；Wen, 2002c），用以描述无能隙量子态中的新型序。理解量子序的一个方式，是看它如何纳入一个一般的序的分类方案（见图 1.1）。首先，不同的序可以分成两大类：对称性破缺序与非对称性破缺序。对称性破缺序可以用局域序参量描述，可以说其中包含点状物体的凝聚，凝聚的振幅对应于序参量。所有对称性破缺序都可以用朗道对称性破缺理论来理解。非对称性破缺序既不能用对称性破缺描述，也不能用与之相关的局域序参量和长程关联描述，因此是一种新型的序。如果一个量子系统（零温下的态）包含非对称性破缺序，那么就说该系统包含非平庸的量子序。我们看到，量子序无非是量子系统中的非对称性破缺序。

量子序还可以进一步分成许多子类。如果一个量子态有能隙，那么相应的量子序就称为拓扑序；拓扑有序态的低能有效理论将是拓扑场论（Witten, 1989）。第二类量子序出现在费米液体（或自由费米子系统）中，费米液体中不同的量子序由费米面拓扑来分类（Lifshitz, 1960）。第三类量子序来自弦网的凝聚（或简称弦网凝聚）（Wen, 2003a; Levin and Wen, 2003; Wen, 2003b），这一类量子序与“粒子”凝聚的对称性破缺序有某些相似之处。

我们知道，不同的对称性破缺序可以用对称群来分类：利用群论，我们可以把三维全部 230 种晶体序分类。对称性还在对称性破缺基态之上产生并保护无能隙集体激发——Nambu–Goldstone 玻色子。类似地，不同的弦网凝聚（及相应的量子序）可以用一个叫做投影对称群（projective symmetry group，PSG）的数学对象来分类（Wen, 2002c）。利用 PSG，我们可以把 100 多种具有相同对称性的不同二维自旋液体分类。正如对称群一样，PSG 也能产生并保护无能隙激发；然而，与对称群不同的是，PSG 产生并保护的是无能隙规范玻色子和费米子（Wen, 2002a,c; Wen and Zee, 2002）。正因为如此，我们可以说，光和无质量费米子可以有统一的起源：它们可以从弦网凝聚中演生出来。

按照图 1.1 中各序的分类，本书可以分为两部分。第一部分（第 3–5 章）讨论“粒子”凝聚产生的对称性破缺序。我们发展有效理论，研究来自序参量涨落的无能隙 Nambu–Goldstone 模的物理性质。这一部分讲述“声音的起源”以及其他 Nambu–Goldstone 模，还讲述超流体中旋子之间 $1/r^4$ 偶极相互作用定律的起源。第二部分（第 7–10 章）讨论量子序/拓扑序与弦网凝聚。同样，我们发展有效理论，研究低能集体模的物理性质。不过在这里，集体模来自凝聚弦网的涨落，并给出规范玻色子和费米子。因此，第二部分提供了“光与电子的一种起源”，以及其他规范玻色子和费米子的起源；它还提供了 $1/r^2$ 库仑定律（或更一般地，电磁学定律）的一种起源。

\section{光与费米子的起源}

\begin{keypoints}
\item 弦网凝聚为光与费米子的起源提供了一个答案。它统一了规范相互作用与费米统计。
\end{keypoints}

\begin{figure}[H]
\centering
\includegraphics[width=0.72\textwidth]{fig_1.2.pdf}
\caption*{图 1.2\quad 我们的真空也许是一个充满弦网的态。弦网的涨落给出规范玻色子。弦的端点对应于电子、夸克等。}
\end{figure}

我们过去相信，要在理论中拥有光和费米子，就必须人为引进一个基本的 $U(1)$ 规范场和反对易的费米子场，因为在那个年代，我们还不知道有任何集体模的行为像规范玻色子和费米子。然而，由于过去二十年的进展，我们现在知道如何构造具有演生\emph{解禁闭}规范玻色子和/或费米子的\emph{局域玻色系统}（Foerster et al., 1980; Kalmeyer and Laughlin, 1987; Wen et al., 1989; Read and Sachdev, 1991; Wen, 1991a; Moessner and Sondhi, 2001; Motrunich and Senthil, 2002; Wen, 2002a; Kitaev, 2003; Levin and Wen, 2003）。特别地，人们可以在立方格子上构造丑陋的玻色自旋模型，其低能有效理论却是美丽的量子电动力学（QED）和量子色动力学（QCD），带有演生的光子、电子、夸克和胶子（Wen, 2003b）。

这引出了如下问题：自然界中的光和费米子，究竟是像 $U(1)\times SU(2)\times SU(3)$ 标准模型中那样，来自基本的 $U(1)$ 规范场和反对易场，还是来自我们真空中某种特定的量子序？库仑定律是自然界的基本定律，还是仅仅是一种演生现象？显然，假定光和费米子以及库仑定律都来自我们真空中的量子序，更为自然。从弦网凝聚、量子序与无质量规范/费米子激发之间的联系，我们看到弦网凝聚提供了一条统一光与费米子的途径。我们非常想就光和费米子的三个基本问题提出如下可能的回答。

\textbf{光和费米子是什么？}

光是（任意大小的）凝聚弦网的涨落；费米子是凝聚弦的端点。

\textbf{光和费米子从哪里来？}

光和费米子来自充满空间的弦网的集体运动（见图 1.2）。

\textbf{光和费米子为什么会存在？}

光和费米子之所以存在，是因为我们的真空碰巧具有一种叫做弦网凝聚的性质。

假如我们的真空当初选择的是“粒子”凝聚，那么低能下就将只有 Nambu–Goldstone 玻色子。那样的宇宙会非常乏味。弦网凝聚以及随之而来的光和费米子，则提供了一个有趣得多的宇宙——至少有趣到足以支撑智慧生命去研究光与费米子的起源。

\section{新奇比正确更重要}

\begin{keypoints}
\item 道可道，非常道；名可名，非常名。无名，天地之始；有名，万物之母。\footnote{这些是中国哲学家老子（Lao Zi）在 2500 多年前所写《道德经》的开头四句。原书英文是对它们的宽松直译。“道”兼有“道路”、“法则”、“品行”等含义。《道德经》存在许多差别很大的译本，上网搜索并比较这些不同的译本是一件趣事。下面是结合本书语境的一个译法：“可以表述的物理理论不可能是终极的理论。可以实施的分类不可能分类万物。不可表述的终极理论确实存在，并支配着宇宙的创生。已表述的理论描述我们每天看到的物质。”}
\item 能被言说的不会是新奇的；不能被言说的不会是正确的。
\end{keypoints}

在这篇引言中（以及本书的某些部分），我希望让读者感受到：在理论凝聚态物理中，我们从哪里来、现在站在哪里、正向何处去。我并不打算在这里总结公认的看法；相反，我想表达的是我个人的、有意夸张的、关于凝聚态物理和高能物理中许多基本问题的观点。这些观点和图像未必正确，但我希望它们具有启发性。从物理学历史的经验来看，我们可以放心地假定：现行的物理理论没有一个会是完全正确的。（按照老子（Lao Zi）的说法，能够写下来的理论不可能是永恒的理论，因为它受限于我们借以写下该理论的数学符号。）问题在于判定现行理论在哪个方面错了、该如何修正。这需要大量的想象力和激励。

\section{评注：基本粒子概念的演化}

\begin{keypoints}
\item 随着时间的流逝，基本粒子的地位不断被降级：从万物的砌块，降为（很可能）只是一个低微玻色模型的集体激发。
\end{keypoints}

地球曾被看作宇宙的中心。随着时间的流逝，它的地位降低为不过是宇宙中亿万颗行星之一。看来，基本粒子的概念可能也会有类似的命运。

在人类文明的开端，人们认识到物体可以被分割成越来越小的部分。中国哲学家提出，这种分割可以无限地继续下去，因而不存在基本粒子。希腊哲学家则假定分割不能无限地继续，于是存在终极而不可分割的粒子——万物的砌块。这可能是基本粒子的第一个概念。那些终极粒子被称为 \emph{atomos}。大量科学研究一直致力于寻找这些 \emph{atomos}。

大约在 1900 年，化学家发现所有物质都由几十种不同的粒子构成。人们操之过急地把它们命名为原子（atoms）。电子发现之后，人们认识到基本粒子比原子更小。现在，许多人相信光子、电子、夸克和其他几种粒子是基本粒子。这些粒子由一个叫做 $U(1)\times SU(2)\times SU(3)$ 标准模型的场论描述。

尽管 $U(1)\times SU(2)\times SU(3)$ 标准模型是一个非常成功的理论，现在大多数高能物理学家相信，它并不是终极的万物理论。$U(1)\times SU(2)\times SU(3)$ 标准模型可能是从一个更深层结构中演生出来的有效理论。问题是：标准模型可能从什么样的结构中演生出来？

一个提案是大统一理论，其中 $U(1)\times SU(2)\times SU(3)$ 规范群被提升为 $SU(5)$ 甚至更大的规范群（Georgi and Glashow, 1974）。大统一理论把 $U(1)\times SU(2)\times SU(3)$ 标准模型中的粒子归入非常漂亮而且简单得多的结构。然而，我想指出，我并不认为在大统一理论中光子、电子和其他基本粒子是演生的。在大统一理论中，规范结构和费米统计是基本的，其含义是：想要有规范玻色子和费米子，唯一的办法就是引进矢量规范场和反对易费米子场。因此，为了拥有光子、电子和其他基本粒子，我们必须人为引进相应的规范场和费米子场。所以，规范玻色子和费米子是被人为加进大统一理论的；它们并不是从更简单的结构中演生出来的。

第二个提案是超弦理论（Green et al., 1988; Polchinski, 1998）。某些超弦模型可以给出有效的 $U(1)\times SU(2)\times SU(3)$ 标准模型，外加许多额外的（近）无质量激发。规范玻色子和引力子是演生的，因为超弦理论本身不包含规范场。然而，费米统计不是演生的：电子和夸克来自 $(1+1)$ 维世界面上的反对易费米子场。我们看到，在超弦理论中，规范玻色子和规范结构不是基本的，但费米统计和费米子仍然是基本的。

最近，人们意识到可能存在第三种可能——弦网凝聚。Banks et al.（1977）和 Foerster et al.（1980）首先指出，光可以作为局域玻色模型的低能集体模演生出来。Levin and Wen（2003）指出，即使是三维费米子，也可以作为凝聚弦的端点从局域玻色模型中演生出来。把这两个结果结合起来，我们发现：只要玻色模型具有弦网凝聚，光子、电子、夸克和胶子（或者更精确地说，$U(1)\times SU(2)\times SU(3)$ 标准模型的 QED 和 QCD 部分）就可以从局域玻色模型中演生出来（Wen, 2002a, 2003b）。这一提案的诱人之处在于，规范玻色子和费米子有着统一的起源。在弦网凝聚的图像中，规范结构和费米统计都不是基本的；所有基本粒子都是演生的。

然而，第三个提案也有一个问题：由于手征费米子问题，我们还不知道如何产生标准模型的 $SU(2)$ 部分。自然界有五个深奥之谜，即全同粒子、费米统计、规范结构、手征费米子和引力。弦网凝聚只回答了前三个谜；还剩两个有待解决。

% ==================================================================
% 新增术语（本章新增，供术语表追加）：
%   phonon 声子；roton 旋子；sound wave 声波；dipolar interaction 偶极相互作用；
%   building block 砌块；black box 黑箱；fractional statistics 分数统计；
%   non-abelian statistics 非阿贝尔统计；Nambu–Goldstone mode/boson Nambu–Goldstone 模/玻色子；
%   spontaneous breaking of continuous symmetry 连续对称性自发破缺；
%   translational symmetry breaking 平移对称性破缺；topological field theory 拓扑场论；
%   fault-tolerant quantum computation 容错量子计算；edge states 边缘态；
%   ground-state degeneracy 基态简并；natural selection 自然选择；
%   grand unified theories 大统一理论；superstring theory 超弦理论；graviton 引力子；
%   gluon 胶子；world sheet 世界面；chiral fermion problem 手征费米子问题；
%   identical particles 全同粒子；random systems 无序系统；many-boson system 多玻色子系统；
%   standard model 标准模型；atomos atomos（希腊语“不可分”，原子一词的来源）
% ==================================================================
